\documentclass[10pt,a4paper]{amsart}
\usepackage[margin=19mm]{geometry}
\usepackage{amsmath,amssymb,mathtools}
\usepackage[scr=rsfs,cal=euler]{mathalfa}
\usepackage{xfrac}
\usepackage[unicode,hidelinks,bookmarks=false]{hyperref}
\newcommand{\ie}{i.e.\ }
\newcommand{\cf}{cf.\ }
\newcommand{\resp}{respectively\ }
\newcommand{\Subsection}{\subsection}
\providecommand{\nobreakdash}{\nobreak\hbox{-}}
\setlength{\emergencystretch}{2em}
\multlinegap0pt
\usepackage{bm}
\usepackage{bbm}
\usepackage{dsfont}
\usepackage{xspace}
\usepackage{cancel}
\usepackage{mathtools}
\usepackage{amsthm}
\makeatletter
\@ifundefined{prop}{\newtheorem{prop}{Proposition}}{}
\makeatother
\newcommand\mf{\mathbf}
\newcommand\ml{\mathcal}
\newcommand{\E}{\mathbf{E}}
\newcommand{\R}{\mathbb{R}}  
\newcommand{\bea}{\begin{equation}\begin{aligned}}
\newcommand{\eea}{\end{aligned}\end{equation}}
\newcommand{\mb}{\mathbb}
\newcommand{\rw}{\rightarrow}
\renewcommand{\l}{\left} 
\renewcommand{\r}{\right} 
\title[Towards Quantifying the Hessian Structure of Neural Networks]{Towards Quantifying the Hessian Structure of Neural Networks}
\author{Zhaorui Dong, Yushun Zhang, Jianfeng Yao, Ruoyu Sun}
\date{Source: arXiv:2505.02809v2 (CC BY 4.0)}
\begin{document}
\maketitle

\noindent\emph{Source and licence.} Towards Quantifying the Hessian Structure of Neural Networks. Version arXiv:2505.02809v2 (CC BY 4.0), distributed under the Creative Commons Attribution 4.0 International licence, as recorded in the arXiv licence metadata for this exact version and shown on the abstract page https://arxiv.org/abs/2505.02809v2. Full rights notice: licenses/arxiv-2505.02809-CC-BY-4.0.txt.\par\smallskip

\noindent\textit{Editorial scope.} Counted unit: the Prop whose source label is \texttt{None} in the article (source0.tex lines 1612--1619, source label \texttt{None}), delivered together with the article's own complete original proof of it (source0.tex lines 1620--1700, one \texttt{proof} environment of 81 lines). No mathematics is written, repaired or re-derived: statement and proof are the article's own text line for line. Required context retained uncounted: MP\_NN\_CE\_Hwiwi (lines 1544--1572); MP\_NN\_CE\_Hwiwj (lines 1544--1572). Every definition, display and label of those lines is retained unchanged. Omitted from this package and not redistributed: 1--1543, 1573--1611, 1701--2368. Nothing inside a retained excerpt is omitted or altered: every retained body is delivered exactly as the article's lines are written, so no omission receipt of any kind accompanies this package. Citations: the retained excerpts carry no citation command at all, so no bibliography entry of the article is transcribed and no reference is redistributed. Reference adaptation: none was needed and none was made; every reference inside the delivered unit resolves inside it, so no unresolved reference or citation is produced. Printed slips kept literal: no printed word, symbol, hypothesis or number of the counted statement or of its proof was altered. Layout and class: the article's own class is \texttt{article} with packages including textcomp, stfloats, url, verbatim, graphicx, cite, fullpage, natbib, mathpazo, inputenc, fontenc, hyperref, booktabs, amsfonts; the wrapper is the lane's pinned \texttt{amsart} editorial wrapper at 10pt on A4, which restates the theorem-like environments the retained text needs and adds the article's own macro definitions that the retained text needs (\texttt{\textbackslash E}, \texttt{\textbackslash R}, \texttt{\textbackslash bea}, \texttt{\textbackslash eea}, \texttt{\textbackslash l}, \texttt{\textbackslash mb}, \texttt{\textbackslash mf}, \texttt{\textbackslash ml}, \texttt{\textbackslash r}, \texttt{\textbackslash rw}), with extra wrapper packages (bm, bbm, dsfont, xspace, cancel, mathtools). The delivered numbering is the unit's own. Assets: no retained span contains a figure environment, a graphics-inclusion command, a PDF-page inclusion, a table or a TikZ picture; no image file is copied into this package, the package declares no asset dependency and no image file is redistributed with it. Licence: version arXiv:2505.02809v2 of this article is distributed under the Creative Commons Attribution 4.0 International licence, as recorded in the arXiv licence metadata for this exact version and shown on the abstract page https://arxiv.org/abs/2505.02809v2. A full rights notice is delivered at licenses/arxiv-2505.02809-CC-BY-4.0.txt. Characters: every retained excerpt is pure ASCII, so the delivered engine, XeLaTeX with its default Latin Modern text font, reports no missing character.
    \begin{enumerate}
        \item $\mu_{G_{ii}}$ converges weakly almost surely to a deterministic measure $\mu^G_{11}$, where $s_{\mu^G_{11}}(z)$ is uniquely specified by the functional equation
        \bea\label{MP_NN_CE_Hwiwi}
        s(z)=\frac{1}{\int_{\mb{R}} \frac{t\nu_{1}(dt)}{1+\gamma s(z)t}
 dt -z},\quad \forall z\in \mb{C}^+. 
\eea
Here $\nu_1$ is defined as follows. Let 
$\mf{z}=(z_1,\cdots,z_m)\sim \mathcal{N}_m(0,I_m)$, define random variables
$$
r_k=\frac{\exp(\sigma(\mf{z})^\top v_k)}{\sum_{l=1}^C\exp(\sigma(\mf{z})^\top v_l)},\quad k\in [C],
$$
$$
\xi_C(V)=\mf{1}(z_1>0)\l[\sum_{k=1}^C r_kV_{k1}^2-\l(\sum_{k=1}^C r_k V_{k1}\r)^2\r].
$$
Then $\nu_1$ is given by that for all intervals $\Delta\subset \mb{R}$, $\nu_1(\Delta)=\mf{P}(\xi_C(V)\in \Delta)$.

        
        \item For $i\ne j$, $\mu_{G_{ij}}$ converges weakly almost surely to a deterministic measure $\mu^G_{12}$, where $s_{\mu^G_{12}}(z)$ is uniquely specified by the functional equation
\bea\label{MP_NN_CE_Hwiwj}
        s(z)=\frac{1}{\int_{\mb{R}} \frac{t\nu_{2}(dt)}{1+\gamma s(z)t}
 dt -z},\quad \forall z\in \mb{C}^+. 
\eea
Here $\nu_2$ is defined as follows. Let 
$\mf{z}=(z_1,\cdots,z_m)\sim \mathcal{N}_m(0,I_m)$, define random variables
$$
\eta_C(V)=\mf{1}(z_1>0)\mf{1}(z_2>0)\l[\sum_{k=1}^C r_kV_{k1}V_{k2}-\l(\sum_{k=1}^C r_k V_{k1}\r)\l(\sum_{k=1}^C r_k V_{k2}\r)\r].
$$
Then $\nu_2$ is given by that for all intervals $\Delta\subset \mb{R}$, $\nu_2(\Delta)=\mf{P}(\eta_C(V)\in \Delta)$.
    \end{enumerate}

\begin{prop}
    Suppose that $m\geq 3$ is fixed, and $V\in\R^{C\times m}$ has i.i.d. $\ml{N}(0,\frac{1}{m})$ entries. If as $d,N\rw\infty, \frac
    {d}{N}\rw \gamma\in (0,+\infty)$, then for $i\ne j$, 
    \bea
\lim_{C\rw\infty}\lim_{d,N\rw\infty}\frac{\E\l[\|G_{ii}\|_{\operatorname{F}}^2\r]}{d}&=\frac{2\gamma+1}{4m^2},\\
\lim_{C\rw\infty}\lim_{d,N\rw\infty}\frac{C\E\l[\|G_{ij}\|_{\operatorname{F}}^2\r]}{d}&=\frac{\gamma(m-1)^2}{2^m(m-2)^3m}\l(\sqrt{\frac{m}{m-2}}+1\r)^{m-2}.
        \eea
\end{prop}

\begin{proof}
Let $\overline{V}$ be a deterministic realization of $V$. By expanding \eqref{MP_NN_CE_Hwiwi} and \eqref{MP_NN_CE_Hwiwj} at $z=\infty$, we have almost surely 
\bea
\lim_{d,N\rw\infty}\frac{\|G_{ii}\|_{\operatorname{F}}^2}{d}=\int_\R x^2\mu_{11}^G(dx)=\gamma\E[\xi_C(\overline{V})^2]+(\E[\xi_C(\overline{V})])^2,
\eea
\bea
\lim_{d,N\rw\infty}\frac{\|G_{ij}\|_{\operatorname{F}}^2}{d}=\int_\R x^2\mu_{12}^G(dx)=\gamma\E[\eta_C(\overline{V})^2]+(\E[\eta_C(\overline{V})])^2.
\eea

Then we have
\bea
\lim_{d,N\rw\infty}\frac{\E\l[\|G_{ii}\|_{\operatorname{F}}^2\ |\ V\r]}{d}=\gamma\E[\xi_C(V)^2|V]+(\E[\xi_C(V)|V])^2,
\eea
\bea
\lim_{d,N\rw\infty}\frac{\E\l[\|G_{ij}\|_{\operatorname{F}}^2\ |\ V\r]}{d}=\gamma\E[\eta_C(V)^2|V]+(\E[\eta_C(V)|V])^2.
\eea
Taking expectation both sides we have
\bea\label{eqRHShii}
\lim_{d,N\rw\infty}\frac{\E\l[\|G_{ii}\|_{\operatorname{F}}^2\r]}{d}=\gamma\E[\xi_C(V)^2]+(\E[\xi_C(V)])^2,
\eea
\bea\label{eqRHShij}
\lim_{d,N\rw\infty}\frac{\E\l[\|G_{ij}\|_{\operatorname{F}}^2\r]}{d}=\gamma\E[\eta_C(V)^2]+(\E[\eta_C(V)])^2.
\eea

Write for short that $\xi_C=\xi_C(V),\ \eta_C=\eta_C(V)$. By the strong law of large number, as $C\rw \infty$,
$$
\E\l[\sum_{k=1}^C r_kV_{k1}^2\ \bigg|\ \mf{z}\r]=\E\l[\frac{\frac{1}{C}\sum_{k=1}^C\exp({\sigma(\mf{z})^\top v_k})V_{k1}^2}{\frac{1}{C}\sum_{k=1}^C\exp({\sigma(\mf{z})^\top v_k})}\ \bigg|\ \mf{z}\r]\rw \frac{\E[\exp({\sigma(\mf{z})^\top v_k})V_{k1}^2\ |\ \mf{z}]}{\E[\exp({\sigma(\mf{z})^\top v_k})\ |\ \mf{z}]}=\frac{\sigma(z_1)^2}{m^2}+\frac{1}{m},
$$
$$
\E\l[\l(\sum_{k=1}^C r_kV_{k1}\r)^2\ \bigg|\ \mf{z}\r]=\E\l[\l(\frac{\frac{1}{C}\sum_{k=1}^C\exp({\sigma(\mf{z})^\top v_k})V_{k1}}{\frac{1}{C}\sum_{k=1}^C\exp({\sigma(\mf{z})^\top v_k})}\r)^2\ \bigg|\ \mf{z}\r]\rw \l(\frac{\E[\exp({\sigma(\mf{z})^\top v_k})V_{k1}\ |\ \mf{z}]}{\E[\exp({\sigma(\mf{z})^\top v_k})\ |\ \mf{z}]}\r)^2=\frac{\sigma(z_1)^2}{m^2}.
$$
Therefore 
$$
\E[\xi_C]=\E[\E[\xi_C|\mf{z}]]=\E\l[\frac{\mf{1}(z_1>0)}{m}\r]=\frac{1}{2m}.
$$
Similarly,
$$
\E[\xi_C^2]=\E[\E[\xi_C^2|\mf{z}]]=\E\l[\frac{\mf{1}(z_1>0)}{m^2}\r]=\frac{1}{2m^2}.
$$
Therefore 
$$
\lim_{C\rw\infty}\lim_{d,N\rw\infty}\frac{\|H_{ii}\|_{\operatorname{F}}^2}{d}=\frac{2\gamma+1}{4m^2}.
$$
For the case of $G_{ij}$, by repeating all arguments above, we have
\bea
\eta_C=&\mf{1}(z_1>0)\mf{1}(z_2>0)\l[\sum_{k=1}^C r_kV_{k1}V_{k2}-\l(\sum_{k=1}^C r_k V_{k1}\r)\l(\sum_{k=1}^C r_k V_{k2}\r)\r]\\
=&\mf{1}(z_1>0)\mf{1}(z_2>0)\sum_{k=1}^C\sum_{l=1}^C r_k r_l V_{k1}(V_{k2}-V_{l2})\\
=&\frac{1}{2}\cdot\mf{1}(z_1>0)\mf{1}(z_2>0)\sum_{k=1}^C\sum_{l=1}^C r_k r_l (V_{k1}-V_{l1})(V_{k2}-V_{l2})\\
=&\frac{1}{2}\cdot\mf{1}(z_1>0)\mf{1}(z_2>0) \frac{\sum_{k=1}^C\sum_{l=1}^C \exp(\sigma(\mf{z})^\top v_k)\exp(\sigma(\mf{z})^\top v_l) (V_{k1}-V_{l1})(V_{k2}-V_{l2})}{\l[\sum_{k=1}^C \exp(\sigma(\mf{z})^\top v_k)\r]^2}.
\eea
Let 
$$
h(k,l)=\exp(\sigma(\mf{z})^\top v_k)\exp(\sigma(\mf{z})^\top v_l) (V_{k1}-V_{l1})(V_{k2}-V_{l2}).
$$
Then for $k\ne k'\ne l\ne l'$,
$$
\E[h(k,l)\ |\ \mf{z}]=0,\quad \E[h(k,l)^2\ |\ \mf{z}]<\infty,\quad
\E[h(k,l)h(k',l')\ |\ \mf{z}]=0,
$$
$$
\E[h(k,l)h(k,l')\ |\ \mf{z}]= \l(\frac{\sigma(z_1)^2\sigma(z_2)^2}{m^4}+\frac{\sigma(z_1)^2+\sigma(z_2)^2}{m^3}+\frac{1}{m^2}\r)\exp\l(\frac{3}{m}\|\sigma(\mf{z})\|^2\r).
$$

Then as $C\rw\infty$,
\bea
\E[\sqrt{C}\eta_C]&=\E[\E\sqrt{C}\eta_C|\mf{z}]]\\
&=\E\l[\frac{1}{2}\mf{1}(z_1>0)\mf{1}(z_2>0)\E\l[\frac{\frac{1}{C^{3/2}}\sum_{k=1}^C\sum_{l=1}^C h(k,l)}{\l[\frac{1}{C}\sum_{k=1}^C \exp(\sigma(\mf{z})^\top v_k)\r]^2}\bigg|\mf{z}\r]\r]\\
&=\E\l[\frac{1}{2}\mf{1}(z_1>0)\mf{1}(z_2>0)\frac{\frac{1}{C^{3/2}}\sum_{k=1}^C\sum_{l=1}^C\E\l[ h(k,l)\bigg|\mf{z}\r]}{\E\l[ \exp(\sigma(\mf{z})^\top v_k)\bigg|\mf{z}\r]^2}\r]+o(1)\\
&=o(1),
\eea

\bea
\E[C\eta_C^2]&=\E[\E[C\eta_C^2|\mf{z}]]\\
&=\E\l[\frac{1}{4}\mf{1}(z_1>0)\mf{1}(z_2>0)\E\l[\frac{\frac{1}{C^3}\l[\sum_{k=1}^C\sum_{l=1}^C h(k,l)\r]^2}{\l[\frac{1}{C}\sum_{k=1}^C \exp(\sigma(\mf{z})^\top v_k)\r]^4}\bigg|\mf{z}\r]\r]\\
&\rw \E\l[\frac{1}{4}\mf{1}(z_1>0)\mf{1}(z_2>0)\frac{4\E\l[ h(k,l)h(k,l')\bigg|\mf{z}\r]}{\E\l[ \exp(\sigma(\mf{z})^\top v_k)\bigg|\mf{z}\r]^4}\r]\\
&=\E\l[\mf{1}(z_1>0)\mf{1}(z_2>0)\l(\frac{\sigma(z_1)^2\sigma(z_2)^2}{m^4}+\frac{\sigma(z_1)^2+\sigma(z_2)^2}{m^3}+\frac{1}{m^2}\r)\exp\l(\frac{1}{m}\|\sigma(\mf{z})\|^2\r)\r]\\
&=\l(\E\l[\l(\frac{\sigma(z_1)^2}{m^2}+\frac{1}{m}\r)\exp\l(\frac{\sigma(z_1)^2}{m}\r)\mf{1}(z_1>0)\r]\r)^2 \l(\E\l[\exp\l(\frac{\sigma(z_1)^2}{m}\r)\r]\r)^{m-2}\\
&=\frac{(m-1)^2}{2^m(m-2)^3m}\l(\sqrt{\frac{m}{m-2}}+1\r)^{m-2}.
\eea
Then from \eqref{eqRHShij} we finish the proof.
\end{proof}

\end{document}
